% Generated draft: requires mathematical and source review.
\documentclass[11pt]{article}
\usepackage[margin=25mm]{geometry}
\usepackage{fontspec}
\setmainfont{Latin Modern Roman}
\usepackage{amsmath,amssymb,mathtools,unicode-math}
\setmathfont{Latin Modern Math}
\usepackage{enumitem,needspace,xurl,hyperref}
\setlength{\parindent}{0pt}
\setlength{\parskip}{5pt}
\setlength{\emergencystretch}{3em}
\newcommand{\N}{\mathbb{N}}
\newcommand{\Z}{\mathbb{Z}}
\newcommand{\Q}{\mathbb{Q}}
\newcommand{\R}{\mathbb{R}}
\newcommand{\C}{\mathbb{C}}
\newcommand{\F}{\mathbb{F}}
\newcommand{\A}{\mathbb{A}}
\newcommand{\PP}{\mathbb{P}}
\newcommand{\E}{\mathbb{E}}
\newcommand{\eps}{\varepsilon}
\newcommand{\dd}{\,\mathrm{d}}
\begin{document}
\section*{UnsolvedMath ID: 700972}
\textbf{Source problem code:} catalogue entry
\subsection{Definitions and mathematical statement}
A finite \(p\)-group is a group of order \(p^n\), where \(p\) is prime and \(n\) is a nonnegative integer. The source does not specify whether the trivial group, corresponding to \(n=0\), is intended to be included.

To make the Frobenius terminology explicit, use the following conventional interpretation, subject to confirmation against the source's omitted cross-references. A Frobenius pair consists of a group \(G\) and a nontrivial proper subgroup \(H<G\) such that
\[
H\cap gHg^{-1}=\{1\}\qquad\text{for every }g\in G\setminus H.
\]
The subgroup \(H\) is a Frobenius complement. Equivalently, the natural transitive action on \(G/H\) has nontrivial point stabilizers and every nonidentity element fixes at most one point. Define the associated kernel set by
\[
K(G,H)=\{1\}\cup\left(G\setminus\bigcup_{g\in G}gHg^{-1}\right).
\]
When this set is a subgroup, it is the Frobenius kernel. For an infinite group, its being a subgroup is an explicit requirement here, rather than an asserted consequence of the Frobenius-pair condition. The ambient Frobenius groups are not required to be finite by the supplied statement.

\begin{enumerate}
\item[(a)] Is there an integer \(p_0\geq 2\) such that, for every prime \(p\geq p_0\) and every finite \(p\)-group \(P\), there is a Frobenius pair \((G,H)\) with \(H\cong P\)? Thus the threshold is uniform over all the finite \(p\)-groups under consideration, with no stated bound on their orders. Under the nontrivial-complement convention above, the source's treatment of the trivial \(p\)-group needs clarification; restricting this quantifier to \(P\ne\{1\}\) would be an additional interpretation requiring review.
\item[(b)] Is it true that, for every group \(A\), there is a Frobenius pair \((G,H)\) for which \(K(G,H)\) is a subgroup and there is an injective group homomorphism
\[
\iota:A\longrightarrow K(G,H)?
\]
No finiteness, countability, or other restriction on \(A\) is supplied.
\end{enumerate}

The supplied record places part (b) in the Archive of Solved Problems and reports an affirmative answer attributed to V. V. Bludov. It labels the overall record partially solved and retains part (a) in the main section. These are supplied source assessments, not independently verified mathematical conclusions. The definitions above require comparison with the unavailable statements of Problem 6.55 and Archive Problem 6.53 before this draft is accepted.
\subsection{Short English statement}
Part (a) asks whether, for every sufficiently large prime \(p\), every finite \(p\)-group occurs as a Frobenius complement, with a threshold independent of the group. Part (b) asks whether every group embeds in a Frobenius kernel; the supplied archive reports an affirmative answer. The precise infinite-group conventions and the treatment of the trivial \(p\)-group require review.
\subsection{Sources}
\begin{itemize}
\item Supplied local UnsolvedMath registry record, permanent ID \(700972\), exact source code \texttt{catalogue entry}, titled \emph{Kourovka Notebook Problem 13.54}. The record attributes the problem to A. I. Sozutov, proposed in 1995.
\item \emph{The Kourovka Notebook (Unsolved Problems in Group Theory)}, 13th issue (1995), Problem 13.54; the supplied background identifies an update in the 21st issue (2026), arXiv:1401.0300v48. The supplied collection description names E. I. Khukhro and V. D. Mazurov as editors.
\item Source cross-references: Problem 6.55 for part (a), and Archive Problem 6.53 for part (b). Their contents were not supplied.
\item Supplied Archive of Solved Problems attribution for part (b): V. V. Bludov, \emph{Siberian Math. J.} \textbf{38}, no. 6 (1997), 1054--1056. The supplied archive reports an affirmative answer; no article title or further bibliographic details are supplied.
\item Status review dated \texttt{2026-10-16}: the supplied canonical label is \texttt{partially\_solved}, with evidence labelled \texttt{published}; its explanation assigns the solved part to the archive and the remaining part to the main section. This is supplied dated curation, not an independently verified literature assessment.
\end{itemize}
\end{document}
